\documentclass[10pt,a4paper]{amsart}
\usepackage[margin=19mm]{geometry}
\usepackage{amsmath,amssymb,mathtools}
\usepackage[scr=rsfs,cal=euler]{mathalfa}
\usepackage{xfrac}
\usepackage[unicode,hidelinks,bookmarks=false]{hyperref}
\newcommand{\ie}{i.e.\ }
\newcommand{\cf}{cf.\ }
\newcommand{\resp}{respectively\ }
\newcommand{\Subsection}{\subsection}
\providecommand{\nobreakdash}{\nobreak\hbox{-}}
\setlength{\emergencystretch}{2em}
\multlinegap0pt
\usepackage{mathtools}
\usepackage{bm}
\usepackage{amssymb}
\usepackage[shortlabels]{enumitem}
\usepackage{xcolor}
\newtheorem{thm}{Theorem}
\newtheorem{lem}{Lemma}
\usepackage{cleveref}
\crefname{thm}{Theorem}{Theorems}
\crefname{lem}{Lemma}{Lemmas}
\newcommand{\CB}{\mathcal{B}}
\newcommand{\CC}{\mathcal{C}}
\newcommand{\CN}{\mathcal{N}}
\newcommand{\CS}{\mathcal{S}}
\newcommand{\CV}{\mathcal{V}}
\renewcommand{\P}{\mathbb{P}}
\newcommand{\R}{\mathbb{R}}
\newcommand{\Z}{\mathbb{Z}}
\renewcommand{\geq}{\geqslant}
\renewcommand{\l}{\ell}
\renewcommand{\le}{\leqslant}
\renewcommand{\leq}{\leqslant}
\title{Random multiplicative functions with periodic weights}
\author{Jeremy Schlitt}
\date{Source: arXiv:2608.00184v1 (CC BY 4.0)}
\begin{document}
\maketitle

\noindent\emph{Source and licence.} Random multiplicative functions with periodic weights. Version arXiv:2608.00184v1 (CC BY 4.0), distributed by arXiv under the Creative Commons Attribution 4.0 International licence, http://creativecommons.org/licenses/by/4.0/, as recorded in the arXiv licence metadata for this exact version and shown on the abstract page https://arxiv.org/abs/2608.00184v1. Full licence notice: \texttt{licenses/arxiv-2608.00184-CC-BY-4.0.txt}.\par\smallskip

\noindent\textit{Editorial scope.} Counted unit: the article's lem A (source0.tex lines 1035--1043). The article's complete original proof of it is delivered with it (source0.tex lines 1044--1112, one \texttt{proof} environment of 69 lines). No mathematics is written, repaired or re-derived: the statement and the proof are the article's own lines, in the article's own notation, and no retained line is substituted or re-flowed.  Retained uncounted as the source-local context the proof needs: lines 181--183; lines 200--206; lines 209--211; lines 232--235; lines 836--841; lines 844--846; lines 851--854; lines 882--888.  Nothing was omitted from any retained span, so the package carries no omission record.  No retained line contains a citation, so the wrapper carries no bibliography and no bibliography entry.  No retained line contains a figure environment, an image-inclusion command or drawing code, so no image is delivered and the package declares no asset dependency.  Editorial formatting only, in addition: an AMS \texttt{amsart} preamble that restates the article's own declarations and macros used by the retained lines, namely the environments \texttt{thm}, \texttt{lem} on separate counters, together with the standard packages those lines need.  The retained environments print in this document as Theorem 1, Lemma 1 in order of appearance, whereas the article prints the same items with the numbering of its own full document.  The complete native source of the article as distributed in the arXiv e-print for this exact version is bundled unchanged as \texttt{sources/206521/source0.tex}.
	\begin{equation}\label{eq:sound-xu badness approx}
		\|q\theta\|:= \min_{n \in \Z}|q \theta -n |\geq C\exp(-|q|^{1/127}) \quad \forall q \in \Z \setminus \{0\}.
	\end{equation}

	\begin{thm}\label{thm:A}
		Fix $\theta$ to be an irrational number satisfying \eqref{eq:sound-xu badness approx}. Let $g \in \CB\CV[\R/\Z]$ be a non-constant function. Then the following two statements are equivalent.
		\begin{enumerate}
			\item  $\{S_{g}(N)\}_{N \geq 1}$ converges in distribution to a complex normal distribution with mean $0$ and variance $1-|\hat{g}(0)|^2/\|g\|_2^2$ as $N \to \infty$
			\item For all positive integers $a,b$, such that $a \ne b$ and $(a,b) = 1$, one has $$\int g(at)\overline{g(bt)}dt = \bigg|\int_0^1 g(t)dt \bigg|^2.$$
		\end{enumerate}
	\end{thm}

		\begin{equation}\label{eq: def Delta}
			\Delta(g) :=  \sum_{\substack{\l_1,\l_2,\l_3,\l_4 \in \Z \\ \l_1\l_2 = \l_3\l_4 \\ \l_1 \ne \l_3, \l_2 \ne \l_4 \\ \l_1 \l_3,\: \l_2\l_4 > 0}}\frac{\hat{g}(\l_1)\hat{g}(\l_2) \overline{\hat{g}(\l_3)\hat{g}(\l_4)}}{H(\l_1/\l_3)H(\l_2/\l_4)},
		\end{equation}

	\begin{thm}\label{thm:B}
		Fix $\theta$ to be an irrational number satisfying \eqref{eq:sound-xu badness approx}. Let $g \in \CB\CV_0[\R/\Z]$ be a non-constant function. Let $\xi(N)$ be a function of $N$ which tends to $+\infty$, and satisfies $N-\xi(N) \to +\infty.$ Then $S_{g}(N;\xi(N))$ converges in distribution to a standard complex Gaussian random variable as $N \to \infty$. That is
		$$S_{g}(N;\xi(N)) \xrightarrow[(N \to \infty)]{d} \CC\CN(0,1).$$
	\end{thm}

	\begin{equation}\label{eq: def z}
		z := \begin{cases}
			1, \quad & \text{ in the context of \Cref{thm:A}},\\
			\xi(N), \quad & \text{ in the context of \Cref{thm:B}}
		\end{cases}.
	\end{equation}

	\begin{equation}\label{eq: def trig poly}
		Q(t) = Q_{L,g}(t) := \sum_{|\l| \leq L} c_\l e(\l t), \quad (c_\l \ll_g |\l|^{-1}, c_0 = 0).
	\end{equation}

	\begin{lem}\label{lem: verif cond 3}
		Let $A > 0$ be given, let $\CS$ be as in \eqref{eq: def z} and let $Q$ be as in \eqref{eq: def trig poly}. Uniformly for $L \geq 1$, one has
		$$\bigg|\sum_{\substack{n_1,\dots,n_4 \in \CS \\ n_1n_2=n_3n_4 \\ P^+(n_1)=P^+(n_2) = P^+(n_3) = P^+(n_4)}}Q(\theta n_1)Q(\theta n_2)\overline{Q(\theta n_3)Q(\theta n_4)}\bigg| \ll_{A,g} \frac{N^2 (\log L)^4}{(\log N)^A}.$$
	\end{lem}

	\begin{lem}\label{lem:reduction cond 2} Let $N$ be large. Fix a number $\theta$ which satisfies \eqref{eq:sound-xu badness approx}. Let $\xi(N)$ be a function of $N$ which tends to $+\infty$ as $N \to \infty$. Let $\CS$ be defined as in \eqref{eq: def z}. Suppose $Q$ is a trigonometric polynomial of the form \eqref{eq: def trig poly}. Then, uniformly for $L \leq \min\{\xi(N),(\log N)^{1/10}\}$, one has
		
		\begin{align*}
			\sum_{\substack{n_1,\dots,n_4 \in \CS \\ n_1n_2=n_3n_4 \\ n_1 \ne n_3,n_2 \ne n_4 \\ P^+(n_1)=P^+(n_3) \\ P^+(n_2) = P^+(n_4)}}Q(\theta n_1)Q(\theta n_2)\overline{Q(\theta n_3)Q(\theta n_4)} &= \sum_{\substack{\l_1,\dots,\l_4 \\ |\l_i| \leq L}}C(\bm{\l}) \sum_{\substack{n_1,\dots,n_4 \in \CS \\ n_1n_2=n_3n_4 \\ \{n_1\l_1,n_2\l_2\}=\{n_3\l_3,n_4\l_4\} \\ n_1 \ne n_3,n_2 \ne n_4 \\ P^+(n_1)=P^+(n_3) \\ P^+(n_2) = P^+(n_4) }} 1\\& +O\bigg(\frac{N^2L}{(\log N)^{ \frac{4}{5}}}\bigg).
		\end{align*}
		where $C(\bm{\l}) := c_{\l_1}c_{\l_2}\overline{c_{\l_3}c_{\l_4}}$. 
	\end{lem}

	\begin{lem}\label{lem:verif cond 2 thm A} Let $N$ be large in terms of $\theta$. Suppose $Q_L$ is a trigonometric polynomial of the form \eqref{eq: def trig poly}, and let $z \in \{1, \xi(N)\}$. Then, uniformly for $L \leq (\log N)^{1/10}$, one has
		\begin{align*}
			\sum_{\substack{n_1,\dots,n_4 \leq N \\ P^-(n_i) > z \\ n_1n_2=n_3n_4 \\ n_1 \ne n_3,n_2 \ne n_4 \\ P^+(n_1)=P^+(n_3) \\ P^+(n_2) = P^+(n_4)}} \hspace{-10pt} Q(\theta n_1)Q(\theta n_2)\overline{Q(\theta n_3)Q(\theta n_4)} = \begin{cases}
				N^2 \Delta(Q_L) + O\bigg(\frac{N^2L}{(\log N)^{\frac{4}{5}}}\bigg) & \text{if } z=1, \\
				O\bigg(\frac{N^2L}{(\log N)^{\frac{4}{5}}}\bigg) & \text{if } z=\xi(N).
			\end{cases}
		\end{align*}
		where $\Delta$ is defined in \eqref{eq: def Delta}.
	\end{lem}

	\begin{proof}
		First of all, by the definition of $Q$ and \Cref{lem:reduction cond 2}, we have
		\begin{equation}\label{eq: verif cond 2 first}
			\sum_{\substack{n_1,\dots,n_4 \in \CS \\ n_1n_2=n_3n_4 \\ n_1 \ne n_3,n_2 \ne n_4 \\ P^+(n_1)=P^+(n_3) \\ P^+(n_2) = P^+(n_4)}}Q(\theta n_1)Q(\theta n_2)\overline{Q(\theta n_3)Q(\theta n_4)}= \sum_{\substack{\l_1,\dots,\l_4 \\ |\l_i| \leq L}} C(\bm{\l}) \sum_{\substack{n_1,\dots,n_4 \in \CS \\ n_1n_2=n_3n_4 \\ \{n_1\l_1,n_2\l_2\}=\{n_3\l_3,n_4\l_4\} \\ n_1 \ne n_3,n_2 \ne n_4 \\ P^+(n_1)=P^+(n_3) \\ P^+(n_2) = P^+(n_4) }} \hspace{-1cm} 1 +O\bigg(\frac{N^2 L}{(\log N)^{\frac{4}{5}}}\bigg).
		\end{equation}
		We will simplify the inner sum on the right hand side, and then sum over $\l$. By an application of Rankin's trick, one can show that for any $A>0$,
		\begin{equation}\label{eq: replace Pn Pnl}
			\sum_{\substack{n_1,\dots,n_4 \leq N \\ n_1n_2 = n_3n_4 \\ \exists i : P^+(n_i) \leq P^+(\l_i)}} 1 \ll \sum_{\substack{n_1,n_2 \leq N \\ P^+(n_1) \leq L}} \tau(n_1n_2) \ll_A \frac{N^2}{(\log N)^A}.
		\end{equation}
		As such, we may assume throughout the course of our calculations that $P^+(n_i) = P^+(n_i\l_i)$ for $1 \leq i \leq 4$, at the cost of an error term $O(N^2/(\log N)^A)$. In particular, we have
		\begin{equation}\label{eq:intermed 1} 
			\sum_{\substack{n_1,\dots,n_4 \in \CS \\ n_1n_2=n_3n_4 \\ \{n_1\l_1,n_2\l_2\}=\{n_3\l_3,n_4\l_4\} \\ n_1 \ne n_3,n_2 \ne n_4 \\ P^+(n_1)=P^+(n_3) \\ P^+(n_2) = P^+(n_4)}} 1 = \sum_{\substack{n_1,\dots,n_4 \in \CS \\ n_1n_2=n_3n_4 \\ \{n_1\l_1,n_2\l_2\}=\{n_3\l_3,n_4\l_4\} \\ n_1 \ne n_3,n_2 \ne n_4 \\ P^+(n_1)=P^+(n_1 \l_1)=P^+(n_3 \l_3) = P^+(n_3) \\ P^+(n_2) = P^+(n_2 \l_2) = P^+(n_4 \l_4) = P^+(n_4) }} \hspace{-1cm} 1 + O_A\Big(\frac{N^2}{(\log N)^A}\Big).
		\end{equation}
		We decompose the sum on the right hand side of \eqref{eq:intermed 1} into two parts, and subtract a double-counting:
		\begin{equation}\label{eq: S1 S2 decomp}
			\sum_{\substack{n_1,\dots,n_4 \in \CS \\ n_1n_2=n_3n_4 \\ \{n_1\l_1,n_2\l_2\}=\{n_3\l_3,n_4\l_4\} \\ n_1 \ne n_3,n_2 \ne n_4 \\ P^+(n_1)=P^+(n_1 \l_1)=P^+(n_3 \l_3) = P^+(n_3) \\ P^+(n_2) = P^+(n_2 \l_2) = P^+(n_4 \l_4) = P^+(n_4) }} 1 = S_1 + S_2 - S_3,
		\end{equation}
		where
		\begin{align*}
			S_1 &:= \sum_{\substack{n_1,\dots,n_4 \in \CS \\ n_1n_2=n_3n_4 \\ n_1\l_1=n_3\l_3, n_2\l_2 = n_4\l_4 \\ n_1 \ne n_3,n_2 \ne n_4 \\ P^+(n_1) = P^+(n_1 \l_1)=P^+(n_3 \l_3) = P^+(n_3) \\P^+(n_2) =  P^+(n_2 \l_2) = P^+(n_4 \l_4) = P^+(n_4)}} 1; \quad S_2 := \sum_{\substack{n_1,\dots,n_4 \in \CS \\ n_1n_2=n_3n_4 \\ n_1\l_1=n_4\l_4, n_2\l_2 = n_3\l_3 \\ n_1 \ne n_3,n_2 \ne n_4 \\ P^+(n_1) = P^+(n_1 \l_1)=P^+(n_3 \l_3) = P^+(n_3) \\P^+(n_2) =  P^+(n_2 \l_2) = P^+(n_4 \l_4) = P^+(n_4)}} 1. \\ \\
			S_3 &:= \sum_{\substack{n_1,\dots,n_4 \in \CS \\ n_1n_2=n_3n_4 \\ n_1\l_1=n_3\l_3= n_2\l_2 = n_4\l_4 \\ n_1 \ne n_3,n_2 \ne n_4 \\ P^+(n_1) = P^+(n_1 \l_1)=P^+(n_3 \l_3) = P^+(n_3) \\P^+(n_2) =  P^+(n_2 \l_2) = P^+(n_4 \l_4) = P^+(n_4)}} 1.
		\end{align*}
		
		The sum $S_3$ is negligible compared to $S_1,S_2$: By omitting many conditions from the $S_3$ sum, and merely using the fact that fixing one of the $n_i$ gives at most $2L$ choices for the other $n_j$, one has
		\begin{align}\label{eq:new S3 bound}
			S_3 \ll \sum_{\substack{n_1,\dots,n_4 \in \CS \\ n_1\l_1=n_3\l_3= n_2\l_2 = n_4\l_4}} 1 \ll N \cdot L^4.
		\end{align}
		
		
		The conditions on $n_i$ in the sum $S_2$ imply that $P^+(n_1) = P^+(n_2) = P^+(n_3) = P^+(n_4)$. Thus, by \Cref{lem: verif cond 3}, we know that
		\begin{equation}\label{eq: S2 bound}
			S_2 \ll_A \frac{N^2}{(\log N)^A}.
		\end{equation}
		It remains to simplify $S_1$. For a given $(\l_1,\dots,\l_4)$, the numbers $(n_1,\dots,n_4)$ being summed in $S_1$ must satisfy $n_1n_2=n_3n_4, n_1 \ne n_3, n_2 \ne n_4$. Because of the conditions $n_1\l_1 = n_3\l_3, n_2\l_2 = n_4\l_4,$ this is equivalent to $\l_1\l_2=\l_3\l_4, \l_1 \ne \l_3, \l_2 \ne \l_4$. Thus, we can rewrite $S_1$ as
		\begin{align*}
			S_1 = \bm{1}_{\substack{\l_1\l_2=\l_3\l_4 \\ \l_1 \ne \l_3, \l_2 \ne \l_4}} \times \sum_{\substack{n_1,\dots,n_4 \in \CS  \\ n_1\l_1=n_3\l_3, n_2\l_2 = n_4\l_4\\ P^+(n_1) = P^+(n_1 \l_1)=P^+(n_3 \l_3) = P^+(n_3) \\P^+(n_2) =  P^+(n_2 \l_2) = P^+(n_4 \l_4) = P^+(n_4)}} \hspace{-1cm} 1.
		\end{align*}
		In the above sum, $n_1,n_3$ are decoupled from $n_2,n_4$, and so we can split the sum into a product:
		\begin{align}\label{eq: S1 halfway reduction}
			S_1 = \bm{1}_{\substack{\l_1\l_2=\l_3\l_4 \\ \l_1 \ne \l_3, \l_2 \ne \l_4}} \times \bigg( \sum_{\substack{n_1,n_3 \in \CS \\ n_1\l_1=n_3\l_3  \\ P^+(n_1) = P^+(n_1 \l_1)=P^+(n_3 \l_3) = P^+(n_3)}} 1\bigg)\bigg( \sum_{\substack{n_2,n_4 \in \CS \\ n_2\l_2 = n_4\l_4 \\ P^+(n_2) = P^+(n_2 \l_2)=P^+(n_4 \l_4) = P^+(n_4)}} 1\bigg).
		\end{align}
		We note that
		\begin{equation} \label{eq: S1 prime factor reduction}
			\sum_{\substack{n_1,n_3 \in \CS \\ n_1\l_1=n_3\l_3  \\ P^+(n_1) = P^+(n_1 \l_1)=P^+(n_3 \l_3) = P^+(n_3)}} 1 = \sum_{\substack{n_1,n_3 \in \CS \\ n_1\l_1=n_3\l_3  \\ P^+(n_i) > |\l_i| }} 1 = \sum_{\substack{n_1,n_3 \in \CS \\ n_1\l_1=n_3\l_3}} 1 + O\bigg(\sum_{\substack{n_1,n_3 \leq N \\ n_1\l_1=n_3\l_3  \\ P^+(n_i) \leq L }} 1\bigg).
		\end{equation}
		By another application of Rankin's trick, one can compute that 
		\begin{equation}\label{eq: proof lem 4.3 1}
			\sum_{\substack{n_1,n_3 \le N \\ n_1\l_1=n_3\l_3  \\ P^+(n_i) \leq L }} 1 \ll \sum_{\substack{n_1\leq N \\ P^+(n_i) \leq L }} 1 \ll_A \frac{N}{(\log N)^A}.
		\end{equation}
		Since $n_i \geq 1$, the equation $n_1\l_1=n_3\l_3$ has no solutions if $\l_1\l_3 < 0$. Let us now separate the two cases for $z$. 
		
		If $z = \xi(N)$, then we are assuming $L \leq \xi(N) = z$. Thus, for any $n_1$ satisfying $P^-(n_1) > z$, we strictly have $P^-(n_1) > |\l_3|$. This guarantees that $(n_1,\l_3) = 1$. The equation $n_1\l_1=n_3\l_3$ then forces $n_1 = n_3,$ which implies $\l_1 = \l_3$. However, this contradicts our summation condition that $\l_1 \ne \l_3$. Therefore, when $z = \xi(N)$, the sum is completely empty, and $S_1 = 0$.
		
		On the other hand, if $z = 1$, the roughness condition is trivial ($\CS = [2,N]$). The equation $n_1\l_1 = n_3 \l_3$ is equivalent to $n_1 \l_1' = n_3 \l_3'$, where $\l_1' = \l_1/(\l_1,\l_3)$ and $\l_3' = \l_3/(\l_1,\l_3)$. Solutions to the equation with $n_1,n_3 \leq N$ are thus parameterized as $n_1 = k \l_3', n_3 = k\l_1'$, where $k \leq N/\max\{|\l_1'|,|\l_3'|\}.$ Thus, for $z=1$, we have:
		\begin{equation}\label{eq: proof lem 4.3 2}
			\sum_{\substack{n_1,n_3 \in \CS \\ n_1\l_1=n_3\l_3}} 1 = \bm{1}_{\l_1\l_3 > 0} \cdot \bigg\lfloor \frac{N}{\max\{|\l_1'|,|\l_3'|\}} \bigg\rfloor = \bm{1}_{\l_1\l_3 > 0}  \cdot \frac{N}{H(\l_1/\l_3)} + O(1).
		\end{equation}
		Combining \eqref{eq: S1 prime factor reduction}, \eqref{eq: proof lem 4.3 1} and \eqref{eq: proof lem 4.3 2} back into \eqref{eq: S1 halfway reduction}, we see that for $z=1$:
		\begin{align}\label{eq: S1 bound}
			S_1 = \bm{1}_{\substack{\l_1\l_2=\l_3\l_4 \\ \l_1 \ne \l_3, \l_2 \ne \l_4\\ \l_1\l_3,\l_2\l_4 > 0}} \times \frac{N^2}{H(\l_1/\l_3)H(\l_2/\l_4)} + O\bigg(\frac{N^2}{(\log N)^A}\bigg).
		\end{align}
		Substituting \eqref{eq: S1 bound}, \eqref{eq: S2 bound} and \eqref{eq:new S3 bound} into \eqref{eq: S1 S2 decomp} shows that for $z=1$,
		$$\sum_{\substack{n_1,\dots,n_4 \in \CS \\ n_1n_2=n_3n_4 \\ \{n_1\l_1,n_2\l_2\}=\{n_3\l_3,n_4\l_4\} \\ n_1 \ne n_3,n_2 \ne n_4 \\ P^+(n_1)=P^+(n_3) \\ P^+(n_2) = P^+(n_4)}} 1  = \bm{1}_{\substack{\l_1\l_2=\l_3\l_4 \\ \l_1 \ne \l_3, \l_2 \ne \l_4\\ \l_1\l_3,\l_2\l_4 > 0}} \times \frac{N^2}{H(\l_1/\l_3)H(\l_2/\l_4)} + O\bigg(\frac{N^2}{(\log N)^A} + N \cdot L^4\bigg).$$
		By choosing $A$ to be an appropriately large constant, and substituting this into \eqref{eq: verif cond 2 first}, we get that for $z=1$,
		\begin{equation}\label{eq: verif cond 2 before last}
			\sum_{\substack{n_1,\dots,n_4 \in \CS \\ n_1n_2=n_3n_4 \\ n_1 \ne n_3,n_2 \ne n_4 \\ P^+(n_1)=P^+(n_3) \\ P^+(n_2) = P^+(n_4)}} \hspace{-0.25cm} Q(\theta n_1)Q(\theta n_2)\overline{Q(\theta n_3)Q(\theta n_4)}= N^2 \hspace{-0.25cm} \sum_{\substack{\l_1,\dots,\l_4 \\ |\l_i| \leq L \\ \l_1\l_2=\l_3\l_4 \\ \l_1 \ne \l_3, \l_2 \ne \l_4\\ \l_1\l_3,\l_2\l_4 > 0}}\hspace{-0.25cm}\frac{ C(\bm{\l})}{H(\l_1/\l_3)H(\l_2/\l_4)} + O\bigg(\frac{N^2 L}{(\log N)^{\frac{4}{5}}}\bigg),
		\end{equation}
		since $N\cdot L^4$ is negligible compared to the other error terms. By using the definition of $\Delta$ from \eqref{eq: def Delta}, we see that the above equation simplifies precisely to $N^2 \Delta(Q_L) +O\big(N^2 L(\log N)^{-4/5}\big)$ when $z = 1$. Combining this with our earlier deduction that the sum is $O\big(N^2 L(\log N)^{-4/5}\big)$ when $z = \xi(N)$ completes the proof of the lemma.
	\end{proof}

\end{document}
